% SOURCE_STATEMENT: "For mirror pairs of Calabi–Yau threefolds over $\\mathbb Q$, find a satisfactory quantum zeta function and quantum $L$-series that retain the essential weight-motive factor and satisfy the expected $T$-duality."
% Preparation and final self-audit: CONTRIBUTING.md prompts C and D; source/audit artifacts retained in the local physics_statements_2026_10_11 batch.
% !TeX program = lualatex
\documentclass[11pt]{article}
\usepackage[margin=25mm]{geometry}
\usepackage{fontspec}
\setmainfont{Latin Modern Roman}
\usepackage{amsmath,amssymb,mathtools}
\usepackage{unicode-math}
\setmathfont{Latin Modern Math}
\usepackage{xurl,hyperref}
\hypersetup{colorlinks=true,urlcolor=blue}
\setcounter{secnumdepth}{0}
\setlength{\parindent}{0pt}
\setlength{\parskip}{5pt}
\setlength{\emergencystretch}{3em}
\newcommand{\sref}[2]{\hyperlink{source:#1:#2}{[S#2]}}
\newcommand{\cataloguescope}[1]{\paragraph{Recorded target.}#1}
\newcommand{\R}{\mathbb R}
\newcommand{\C}{\mathbb C}
\newcommand{\N}{\mathbb N}
\newcommand{\Z}{\mathbb Z}
\newcommand{\Q}{\mathbb Q}
\newcommand{\F}{\mathbb F}
\newcommand{\eps}{\varepsilon}
\begin{document}
\section{Quantum Zeta Functions for Mirror Calabi–Yau Threefolds}\label{statement}

\subsection{Definitions and mathematical statement}

\paragraph{Shared notation.}$\mathbb N=\{1,2,\ldots\}$; the standard number fields have their usual meanings. All additional conventions are stated locally.


\paragraph{Definitions and conventions.} For a smooth Calabi--Yau threefold $X$ over $\mathbb Q$ and a good prime $p$, its local zeta function encodes point counts, $Z_p(X,T)=\exp(\sum_{m\ge1}|X(\mathbb F_{p^m})|T^m/m)$; its cohomological factors are characteristic polynomials of Frobenius. An $L$-series assembles corresponding local factors over primes. A mirror pair exchanges the appropriate Hodge data. The weight motive in the report's explicit orbifold families contributes a common essential factor to the third-cohomology Frobenius polynomials of the pair.

\paragraph{Source scope and qualifications.} Quantum zeta function and quantum L-series are requested notions, not uniquely defined inputs. The report's ratio candidate has the desired duality but discards the weight-motive factor; its explicit orbifold results are context, not a universal construction for all Calabi--Yau mirror pairs.

\paragraph{Catalogue mathematical target.}
\paragraph{Catalogue wording retained for traceability.}
For mirror pairs of Calabi–Yau threefolds over $\mathbb Q$, find a satisfactory quantum zeta function and quantum $L$-series that retain the essential weight-motive factor and satisfy the expected $T$-duality.

\paragraph{Saved literature context.} The supplied catalogue's assessment dated 2026-08-21 records: The report proves restricted motive and slope-zeta statements, while later work still treats the general arithmetic and quantum zeta-function correspondence as conjectural. The assessment is retained as dated provenance; it is not a new comprehensive literature or open-status certification.
\subsection{Short English statement}
Construct arithmetic mirror invariants with the expected duality while retaining the important motive information lost by simple zeta-function ratios.
\subsection{Sources}
\begin{itemize}
\item[\hypertarget{source:statement:1}{S1}] ULAM.AI and the UnsolvedMath Contributors, UnsolvedMath, catalogue entry. Catalogue curation under CC BY 4.0; definitions and English exposition added locally.
\url{https://huggingface.co/datasets/ulamai/UnsolvedMath}.
\item[\hypertarget{source:statement:2}{S2}] Mini-Workshop: Zeta Functions, Index and Twisted K-Theory; Interactions with Physics (2006), Oberwolfach Reports, DOI 10.4171/owr/2006/21, pp. 1245--1284; official report source and mathematical context.
\url{https://doi.org/10.4171/owr/2006/21}.
\item[\hypertarget{source:statement:3}{S3}] Official EMS Press PDF of the same Oberwolfach report; PDF pages 29--31, printed pages 1273--1275.
\url{https://ems.press/content/serial-article-files/46051}.
\item[\hypertarget{source:statement:4}{S4}] Wan, Mirror Symmetry for Zeta Functions; reference in the saved dated catalogue literature review.
\url{https://arxiv.org/abs/math/0411464}.
\item[\hypertarget{source:statement:5}{S5}] Zeta functions of alternate mirror Calabi-Yau families; reference in the saved dated catalogue literature review.
\url{https://arxiv.org/abs/1612.09249}.
\end{itemize}
\label{end:statement}
\end{document}
